%%%%%%%%%%%%%%%%%%%%%%%%%%%%%%%%%%%%%%%%%%%%%%%%%%%%%%%%%%%%%%%%%%%%%%
% Lớp 4 | Chủ đề 1. Ôn tập và bổ sung | Biểu thức chứa chữ
% Dạng 4D1?4-2: Tìm thành phần chưa biết qua nhiều bước
%%%%%%%%%%%%%%%%%%%%%%%%%%%%%%%%%%%%%%%%%%%%%%%%%%%%%%%%%%%%%%%%%%%%%%


% nguồn: 2026-11-18_lop4_chu-de-1-tu-luan (nhập 18/11/2026)
\begin{bt}%[4D1N4-2]%[ma:d473fed4]
% nguon: Toán tư duy lớp 4 – Bài 4 SGK Toán 4 KNTT
	Tìm $x$: $x - 2\,345 = 6\,789$.
	\loigiai{
		$x$ là số bị trừ. Muốn tìm số bị trừ, ta lấy hiệu cộng với số trừ.\\
		$x = 6\,789 + 2\,345 = 9\,134$.
	}
\end{bt}

% nguồn: 2026-11-18_lop4_chu-de-1-tu-luan (nhập 18/11/2026)
\begin{bt}%[4D1H4-2]%[ma:b975ab09]
% nguon: Toán tư duy lớp 4 – Bài 4 SGK Toán 4 KNTT
	Tìm $x$: $x \times 5 + 125 = 1\,000$.
	\loigiai{
		Coi $x \times 5$ là một số hạng: $x \times 5 = 1\,000 - 125 = 875$.\\
		$x = 875 : 5 = 175$.\\
		Thử lại: $175 \times 5 + 125 = 875 + 125 = 1\,000$.
	}
\end{bt}

% nguồn: 2026-11-18_lop4_chu-de-1-tu-luan (nhập 18/11/2026)
\begin{bt}%[4D1V4-2]%[ma:d6e081a3]
% nguon: Toán tư duy lớp 4 – Bài 4 SGK Toán 4 KNTT
	Tìm $x$: $(x + 15) \times 4 - 60 = 200$.
	\loigiai{
		Gỡ dần từ ngoài vào trong.\\
		$(x + 15) \times 4 = 200 + 60 = 260$.\\
		$x + 15 = 260 : 4 = 65$.\\
		$x = 65 - 15 = 50$.
	}
\end{bt}

% nguồn: 2026-11-18_lop4_chu-de-1-tu-luan (nhập 18/11/2026)
\begin{bt}%[4D1C4-2]%[ma:ba20532c]
% nguon: Toán tư duy lớp 4 – Bài 4 SGK Toán 4 KNTT
	Tìm $x$: $x + x \times 2 + x \times 3 = 4\,200$.
	\loigiai{
		Viết $x = x \times 1$: $x \times 1 + x \times 2 + x \times 3 = x \times (1 + 2 + 3) = x \times 6$.\\
		$x \times 6 = 4\,200$ nên $x = 4\,200 : 6 = 700$.
	}
\end{bt}
